\chapter[Feeling Good or Living Well?]{Happiness: Feeling Good or Living Well?}
\section{A Slip of Paper with One Question}
First pick up something very light: a narrow strip of paper.

It is neither a banknote, a ticket, nor a certificate. It is a questionnaire handed out in a psychology lesson, with one sentence and a scale:

``Overall, how satisfied are you with your life at present? Give a score from 0 to 10, where 0 means completely dissatisfied and 10 completely satisfied.''

No problem to solve, no multiple-choice options, no standard answer. Holding your pen, you have only one task: write a number.

Notice what is strange about this object. Other measuring instruments seem more solid: thermometers measure temperature, scales weight, trackside stopwatches speed. Their objects have recognised locations: temperature under your arm, weight beneath your feet, speed on the track. But where does this slip's object live? Which part of your body contains your ``life satisfaction''? Is this morning's value the same as last night's? If you give yourself 7 and your deskmate gives 6, do you really possess one more unit of that thing?

Stranger still, many people take this slip seriously. Governments spend large sums asking thousands of people almost the same question every year. Organisations turn their answers into a World Happiness Report ranking countries. Some states, such as Bhutan in the Himalayan foothills, even make national happiness a policy goal, pursuing not only gross national product but ``Gross National Happiness.'' One sheet, one question, supports an enormous structure.

This chapter weighs that slip: what does it measure, and what does it miss? Dig further and we encounter a question debated for over two millennia: is happiness ``feeling good'' or ``living well''? These similar-sounding expressions lead down entirely different roads.

\section{Seven Minutes in Psychology Class}
Now enter a scene.

Time: Wednesday afternoon's third period, psychology. Place: an eighth-grade classroom. Outside stands an old camphor tree whose leaf shadows sway across desks in the wind. Ms Gu has just distributed a stack of slips by group.

``No lecture today,'' she says. ``An exercise instead. It is anonymous: no names. We will collect these and calculate the class's average satisfaction.''

Paper rustles for seven minutes. More happens in those minutes than in a whole lesson.

Lin Meng, in the third row, writes 9 almost immediately. She seems to have everything: above-average grades, a comfortable family, a new phone, and invariably bright social-media photographs. Yet a second before writing, she wonders how much of that 9 is real. Last night she watched short videos until half past one; putting down her phone, she felt an indescribable emptiness. The thought lasts one second. She still writes 9 because she ``cannot think of a reason to be dissatisfied.''

Across the aisle, Zhou Ye holds his pen suspended for a long time before writing 4. Home has been troubled: his parents argue, and his mother has gone to his grandmother's. Yet last weekend the robotics team he leads won a district prize, and he shouted himself hoarse with joy at the presentation. The 4 and the trophy inhabit the same person; neither yields to the other.

In the back row sits a recently transferred boy everyone calls Old Ma. He turns the slip over twice, writes nothing, folds it, and pockets it. After class, Ms Gu asks whether he understood. His reply leaves her silent for several seconds: ``I understood. I just don't know which me it is asking about.''

The collected slips lie on the teacher's desk. Ms Gu calculates and announces an average of 6.8, ``within the normal range for this district.''

Some laugh; others shrug. A boy behind Lin Meng mutters, ``Can scores like these even be averaged? My 8 today isn't the same as yesterday's 8.'' Someone adds, ``And who tells the truth in psychology class anyway?''

Remember these seven minutes. One question in a classroom of over forty people has produced at least three tangles. Lin Meng's: everything I have seems right, so why does it feel wrong? Zhou Ye's: the worst and best happen in one week, so how do I total the account? Old Ma's is sharpest: if there is more than one ``me,'' who receives the score?

The rest of this chapter tries to untangle these three knots.

\section{What Does Happiness Ask About?}
State the conflict before rushing to an answer.

In ordinary speech, ``happiness'' refers to at least four completely different things that we never distinguish.

First, \textbf{pleasure}: waves of good feeling. The sweetness of the first mouthful of ice cream, the instant you finish a game, a Sunday morning under the duvet with no need to rise. It comes and goes quickly, measured in bursts.

Second, \textbf{satisfaction}: an overall assessment of your life. Rather than a passing mood, it looks back: ``On balance, have things been going well?'' The questionnaire's 0--10 usually asks about this.

Third, \textbf{desire fulfilment}: getting what you want. Admission to your preferred school, the shoes you wanted, your crush liking you back. Its ledger is simple: the more you obtain, the happier you are.

The fourth is hardest to describe; call it \textbf{living well}. Your life has substance: doing what you should, becoming a decent person, living meaningfully. This can coexist with feeling bad. Someone caring for an ill relative may be exhausted and often distressed, yet answer yes when asked whether their life is worthwhile.

Now the real conflict can appear. Imagine two worlds, only one of which you may choose:

World A: everyone feels wonderful all the time, with unlimited pleasure and no worry, pain, or anxiety. But these feelings are manufactured. Nobody actually does or experiences anything; they merely soak in pleasant sensations.

World B: people experience joy and sorrow, failure, heartbreak, and exhaustion. But they genuinely act, love, make mistakes, and take responsibility. Looking back, they can say, ``This is a life I have lived.''

If happiness means feeling good, World A wins outright. Technology's proper destination is a pleasure-supplying machine, and refusing entry is foolish. If happiness means living well, World A is a disaster: nobody there has truly lived even one day, however sweetly they smile.

This is not merely a science-fiction fancy. It is a real fork in the road, with distinguished thinkers lining both sides. One side says that feeling good is everything and suffering the only certain evil. The other says feelings deceive: what matters is a life worthy of a human being. For over two millennia, brilliant minds have occupied both sides.

Your 0--10 slip sits precisely at the fork. Each side reads its question differently.

Before continuing, choose a side yourself: if you must choose, would you enter World A?

\section{State Both Sides' Strongest Cases}
The thinkers on each side deserve individual introductions. First comes the feeling-good camp. Our rule is to state its case at its strongest before hearing the opposition.

\textbf{First line-up: feeling good is what counts.}

State this position fully. Caricaturing it as ``nothing but eating, drinking, and entertainment'' is unfair.

Start with an ancient Greek misunderstood for over two thousand years: Epicurus. An English word for pleasure-seeking derives from his name, yet his own life was almost austere. Teaching in a garden outside Athens, he held pleasure to be the only good, but defined it not as stimulation but as \textbf{absence of pain}: no bodily suffering or mental anxiety. He said that with bread, water, and conversation with friends, he wanted nothing more. This sounds like a quiet study, not a palace of excess. His misunderstood life reminds us that even the feeling camp's most famous representative knew that good feelings are not the same as intense stimulation.

Now consider a more thoroughgoing advocate: the nineteenth-century British reformer Jeremy Bentham. A founder of \textbf{utilitarianism}, he held that actions, laws, and policies should be judged by whether they produce ``the greatest happiness of the greatest number.'' Happiness meant pleasure minus pain. His argument had three increasingly powerful levels.

First: feeling is the unavoidable judge. We can debate a law's justice or a life's nobility for a century, but a person knows whether they hurt or suffer without theoretical permission. Morality grounded in feeling rests on something everyone possesses.

Second: this accounting treats the vulnerable fairly. In an age when poor people scarcely counted as human, Bentham insisted that a beggar's pain and an aristocrat's pain had equal value. Do not underestimate the radical equality of counting everyone's pleasure once, and only once. He extended the principle beyond humanity: membership in the moral ledger depends not on intelligence but on \textbf{the capacity to suffer}.

Third: this ledger makes rulers answerable. A king invokes national glory, a cleric your soul; Bentham opens the account book. Set glory and souls aside: does this policy make actual people happier or more miserable? That question punctures many abuses committed under lofty banners.

This is a powerful case. Its pleasure-minus-pain ledger helped produce modern gains in public health, labour protection, and animal welfare.

\textbf{Second line-up: feelings deceive; people need to live well.}

The opposing line-up is equally distinguished and draws from entirely different civilisations.

Aristotle leads the ancient Greeks. His famous concept eudaimonia is often translated as ``happiness,'' misleadingly: literally a condition of having a good attendant spirit, it actually means \textbf{flourishing throughout life and realising what a human being can be}. The Nicomachean Ethics argues, broadly, that happiness is not a mood but an assessment of a life's activity, requiring full exercise of distinctively human capacities: reason and virtue. A tree's goodness is not whether it feels comfortable now but whether it has grown into a flourishing tree. Likewise, Aristotle would say the person in the pleasure machine has nothing to do with happiness, however sweet every second: they have realised nothing, only been fed.

At roughly the same time, across the world in China, Confucians offered a structurally similar answer with a different flavour. Confucius discussed not abstract ``happiness'' but joy rooted in human relationships: learning, friends arriving, a clear conscience, and doing right by others. For Confucians, a life's quality cannot be calculated behind a closed door. It depends on standing rightly within the network of family, friends, and the world. Purely private happiness receives a discount in their accounts.

Further southwest, ancient Indian Buddhism took a more surprising route. Dukkha, in Pali, is commonly translated as suffering, but means more than pain and sadness: a condition of \textbf{instability that cannot be securely grasped}. Pleasure belongs to it too because pleasure passes; tighter clinging brings greater pain when it goes. The Four Noble Truths teach, broadly, that life contains this insecurity; craving and attachment cause it; it can cease; and there is a path to cessation. Notice the turn: where others ask how to obtain happiness, Buddhism first asks whether insisting ``I must be happy'' is itself part of the problem. Someone who stops desperately demanding happiness from the world may instead find profound peace.

This line-up's latest arrivals are its bluntest. Twentieth-century French \textbf{existentialists}, including Sartre and Camus, say the question ``What is happiness?'' points in the wrong direction. The world owes you no meaning, and meaning is not for sale. Condemned to freedom, people must \textbf{create} meaning through choices and action in a world without instructions. Camus invokes Sisyphus, condemned by the Greek gods to push a rock uphill forever, only for it to roll down again. Others see absolute despair; Camus's point can be paraphrased like this: Sisyphus can master his fate, knowingly returning down the mountain despite the absurdity. That ``despite'' embodies human dignity. Existentialism asks not whether you feel good but whether you dare genuinely to live.

The line-ups are complete. They disagree not over whether happiness is desirable---nobody opposes it---but over what the word means. When one word houses several meanings and people fight while holding different ones, we can take the confusion apart. Here is the tool.

\section[Thinking Bridge: Four Questions]{Thinking Bridge: Turn Happiness into Four Questions}
Next time ``happiness'' appears in a questionnaire, advertisement, your parents' speech, or your own thoughts, do not immediately answer yes or no. Divide it into four questions. This compresses over two thousand years of argument into a pocket-sized card.

\textbf{Question one: how does it feel?} Pleasure and pain now, arriving in bursts and lasting seconds or hours. The best measurement asks during the experience rather than afterwards. Researchers have actually given participants pagers that sound randomly several times a day, prompting immediate reports of feeling.

\textbf{Question two: am I satisfied overall?} A retrospective assessment, asking not about mood but whether you endorse your life. This is the psychology slip's 0--10 question.

\textbf{Question three: did I get what I wanted?} The desire ledger: each gain counts clearly. Its hidden trap is that wants rise alongside possessions, making the account impossible to settle. More on that shortly.

\textbf{Question four: am I living well?} Virtue and meaning: have I done what I should, become a decent person, lived a substantial life? Aristotle and Confucians primarily occupy this level. Existentialists live nearby but return the standard of ``well'' to each individual.

Apply the card to real statements and many arguments reveal themselves immediately.

An advertisement says, ``Own this: you deserve happiness.'' It sells question one, a burst of pleasure, and question three, obtaining a desired object, while suggesting you are buying question two. Parents say, ``Today's restrictions are for your future happiness.'' They mean question two or even four; what you lose now belongs to one. Neither side need be wrong, but without distinguishing levels, repeating the quarrel ten thousand times solves nothing. Ms Gu's 6.8 measures only question two. Lin Meng's emptiness belongs to one, Zhou Ye's trophy to four, and Old Ma's question lies deeper: with more than one ``me,'' whose overall evaluation counts?

The tool is not four terms to memorise but a reflex to cultivate: whenever happiness is mentioned, ask which level is meant. Many apparently profound contradictions turn out to be people occupying different floors and talking past each other.

\section{A Basket of Rice, a Gourd of Water}
Now read a short primary source. Two and a half millennia ago, someone was already answering whether feelings or one's way of life mattered more.

In the Analects chapter ``Yong Ye,'' Confucius praises his favourite student, Yan Hui:

\begin{quote}
The Master said, ``How worthy Hui is! A basket of food, a gourd of water, living in a poor alley: others could not endure the hardship, yet Hui does not change his joy. How worthy he is!''
\end{quote}
In other words: Yan Hui is truly virtuous. With one bamboo basket of rice and a gourd of water in a shabby lane, where others would be overwhelmed by worry, his joy remains unchanged. Truly virtuous!

In ``Shu Er,'' Confucius speaks about himself:

\begin{quote}
The Master said, ``Eating coarse grain, drinking water, and bending my arm for a pillow, I find joy there too. Wealth and honour gained unjustly are to me like passing clouds.''
\end{quote}
That is: even coarse food, cold water, and a bent arm as a pillow contain joy. Riches and status obtained wrongly matter no more to me than clouds in the sky.

Set these beside the psychology slip and their apparent oddity appears. On question one, Yan Hui's material conditions offer almost no pleasure. On question three, a basket of rice and gourd of water would scarcely feature on anyone's wish list. Yet Confucius says ``joy,'' and unchanged joy: steady and unshakeable. It inhabits question four. Yan Hui lives well; he knows it, and so does Confucius.

Two easy misreadings need clearing up. Distinguish what ancient people said from what we assume they said.

First, Confucius is \textbf{not} praising poverty. The crucial phrase is ``wealth and honour gained unjustly.'' He rejects unjust riches, not riches themselves. If wealth's source is wrong, it is as light as clouds; legitimate wealth carries no implied guilt. Reading Confucianism as ``the poorer, the more glorious'' is a later mistake. Using it to tell poor people to stay obedient misapplies it still further.

Second, unchanged joy does not mean inability to suffer. The five Chinese characters rendered ``others could not endure the hardship'' show that these conditions are objectively difficult. Yan Hui can feel hardship, but his joy does not \textbf{depend} on hardship's absence. This distinction matters enormously: the feeling camp says pain diminishes pleasure; Yan Hui presents another structure, where joy rests on something deeper and suffering reduces only the level of feeling.

These passages matter not because Confucius must be right---you may legitimately think Yan Hui was putting on a brave face---but because they document something. The distinction between feeling good and living well is no modern academic word game; it is among humanity's oldest questions about itself. Ancient Chinese and Greeks disputed it, and Indians approached it from another direction. Today's classroom slip is merely the latest examination paper in that long argument.

\section{One Curve, Three Interpretations}
Now add modern evidence and watch three explanations compete.

First, a repeatedly observed fact. Distinguish what happens from why: only this opening part states the facts. Decades of international surveys broadly show a curve relating money and happiness that rises steeply before flattening. Among the very poor, extra income markedly improves well-being: adequate food, medical care, and schooling each remove real suffering. Beyond basic security and modest prosperity, however, the curve slows, with additional money bringing smaller gains. In the 1970s, American economist Richard Easterlin formally identified a puzzle: within a country, richer people generally report greater happiness, yet decades of national enrichment do not bring matching increases in average happiness. This became the ``Easterlin paradox.''

Those are the facts. Why does this happen? Three explanations each hold part of the truth.

\textbf{Explanation one: diminishing returns (economics).} Each additional unit brings less satisfaction: the first steamed bun saves your life; the tenth is unnecessary. Money first solves life-threatening problems, then important ones, then merely nice-to-have improvements. Naturally its marginal contribution shrinks. This explanation is tidy and calculable, but treating flattening as obvious does not answer the paradox's sharper half: why does collective enrichment leave the average almost unchanged? Even if each bun offers less utility, thousands of additional buns for the nation should leave some gain.

\textbf{Explanation two: a moving ruler (psychology).} Two cunning mechanisms matter. ``Hedonic adaptation'' means getting used to almost every good thing: a new home's delight lasts perhaps months before it becomes simply ``home.'' Happiness resembles a treadmill: you run forward while its belt retreats. ``Social comparison'' means judging your situation less by absolute standards than by other people. When a neighbour replaces their car, yours stays still but your satisfaction with it drops. Together these explain the half economics misses: national enrichment moves everyone's reference point upwards, so nobody wins. Yet the explanation has blind spots. If adaptation and comparison consume everything, do better medicine, rising literacy, and fewer wars leave no trace on happiness? That cannot be right.

\textbf{Explanation three: biased measurement.} This approach inspects the ruler before disputing human nature. Nobel laureate and psychologist Daniel Kahneman distinguished two selves: the \textbf{experiencing self}, living minute by minute, and the \textbf{remembering self}, summarising and scoring afterwards. Their accounts often disagree. His ``peak--end rule'' suggests that remembered evaluations depend hardly at all on average feeling throughout an experience, but instead on feelings at \textbf{its peak} and \textbf{its end}. A largely uneventful trip with a spectacular ending can therefore seem happier in memory than a consistently comfortable one. Who holds the pen over the questionnaire? The summarising, storytelling remembering self, shaped by the peak--end rule. The 0--10 scale may measure not life itself but life's compressed memory file. Part of the curve may be the ruler's shape, not the world's. Yet this cannot explain everything: the consistently low scores across areas of extreme poverty and suffering plainly reflect more than ways of filling in forms.

Each explanation works partly and falls short of the whole. This is not evasive compromise but a methodological reminder: facing a curve, ask how much belongs to the world, how much to the mind, and how much to the ruler.

\section[Further Challenge: The Pleasure Machine]{Further Challenge: The Machine Promising Endless Happiness}
Now introduce our weightiest thought experiment. Robert Nozick presented it in his 1974 book Anarchy, State, and Utopia. It became known as the ``experience machine.''

Imagine super-neuroscientists building a machine. Lie down, attach electrodes, and it supplies any life experience you want, in its finest version: entering your dream school, writing a great novel, being universally loved, watching sunrise from a summit. Everything feels completely real; inside you never know you are in a machine. Before entering, you can customise the story. It carries a lifetime guarantee and never malfunctions.

The sole catch: your body remains lying in a tank.

Would you connect?

Nozick reports that most people hesitate and refuse. That hesitation is a gold mine. If happiness is simply feeling good, the machine is perfect: purer, steadier, more plentiful good feelings than real life. Refusal makes no sense. Yet we hesitate. Nozick says this demonstrates concern for at least three things beyond feeling, which we can summarise as follows: \textbf{first, we want actually to do things, not merely experience doing them; second, we want to be real people, not indeterminate beings floating in tanks; third, we want contact with a real world, not a film screened forever for an audience of one.}

Weigh both the experiment's power and its limits.

Consider an objection that cannot be dismissed easily: people refuse because they are afraid, unfamiliar with the machine, or attached to existing relationships, not because their choice is rational. This has force. If safety were certain, family and friends could connect too, and everyone joined, how many reasons to refuse would remain? That makes the experiment harder, not easier.

Now consider an \textbf{easily misjudged counterexample} that prevents overreach. You might conclude that all manufactured good feelings are false and inferior, and all real suffering deserves praise. That inference is dangerously wrong. Imagine someone enduring relentless agony, such as severe pain in the final stage of an incurable illness, and someone else entirely well. If artificial good feelings could greatly relieve the first person's suffering and bring peace to their remaining life, could we say, ``It is unreal and inferior; you should refuse''? Most of us could not. The experiment challenges replacing an entire life with sensations, not the importance of feelings. For someone suffering, relief is enormously important. That is why Bentham's camp still stands. A thought experiment is a searchlight, not a wall: it illuminates ``happiness exceeds feeling,'' not ``feeling is worthless.''

Existentialists would sharpen the objection further. The machine's deepest problem, they might say, is not unreality but that it \textbf{eliminates your possibility of becoming an author}. However splendid its script, someone else wrote it and a programme performs it; you are a spectator fed sweets. Humans, in this account, must write their own plays in an absurd, unscripted world. Sisyphus has dignity because his stone, his futility, and his decision to continue are real. However sweet false happiness tastes, it cannot answer, ``Is this a life I have made?''

Existentialism must accept a reply too: ``real suffering is better than false sweetness'' carries different weight from a secure, well-fed philosopher than from someone actually suffering. Who may refuse the machine on another person's behalf? Keep that question open.

\section[Pocket Machines and National Rankings]{The Pleasure Machine in Your Pocket and Countries in the Rankings}
This two-thousand-year-old question now refreshes in your pocket every fifteen seconds.

Short-video recommendation algorithms are an undeclared prototype of the experience machine. They promise not a lifetime but the next fifteen seconds, precisely calculating your dopamine and turning question-one good feelings into a continuous conveyor belt. Our tool reveals something important: after an hour of scrolling, someone may have gained many little points for feeling, yet shake their head when asked whether that hour was satisfactory. Of four forms of happiness, the app supplies only the cheapest, while quietly raising the stimulation threshold so ordinary pleasures seem slow and weak. It is not the experience machine, but it trains you to adapt to one.

Consumer culture speaks the same language. ``You deserve it'' is desire-fulfilment happiness's standard slogan. It hands you question three's ledger and helpfully keeps adding wants. Hedonic adaptation waits behind it: acquisition marks the peak, decline follows, and another deserved purchase becomes necessary. This is not one unusually wicked brand but a business system built on confusing happiness's four forms, selling first-level goods at second- and fourth-level prices.

Stopping here would fall into another trap: making happiness entirely personal, as though seeing through advertisements, switching off your phone, mastering Confucius and Aristotle, and adjusting your attitude were enough. This popular version particularly needs scrutiny. Happiness is never just self-cultivation; it has \textbf{public conditions}.

Look at the evidence. The annual World Happiness Report asks people across countries a question much like our slip's; Finland, Denmark, and Iceland regularly lead. The report itself does not attribute this to a superior Nordic mindset. It repeatedly discusses concrete public conditions: affordable treatment, protection against immediate ruin after unemployment, broadly trustworthy officials, trust between strangers, and predictable futures. The leaders are not peoples best at ``looking on the bright side,'' but societies doing the public work of making life secure. Conversely, war, destitution, high unemployment, and pervasive distrust consistently depress whole societies' scores. Telling people there that happiness is a matter of attitude is cruel, not wise.

This leads to politics. Bhutan's king proposed Gross National Happiness in the 1970s, arguing that development should pursue more than output. Increasingly, countries measure subjective well-being alongside production. Both sides of the resulting debate deserve a hearing. One says government should provide security, healthcare, education, and fairness, leaving whether citizens feel happy and what happiness means to them. The other asks what output is for if development does not ultimately improve people's lives: governing cannot mean counting money alone. Unresolved, this debate brings the ancient feeling-good versus living-well dispute to the policy table. Whoever defines happiness holds the greatest power there, something every future voter should remember.

You can now write a full assessment of the questionnaire. It measures broad retrospective evaluations of life across a society; averaged over thousands, these reliably distinguish secure prosperity from widespread misery. As a rough public-health-style gauge it is useful, and discarding it would be unwise. It misses question one, countless moments compressed beyond recognition by peak--end memory; the gap between questions three and four, since someone may fulfil every desire yet live emptily, while Yan Hui can possess almost nothing and retain his joy, both reduced to dry numbers; and performance while answering. The pupil's ``Who tells the truth in psychology class?'' names a recognised measurement problem: social desirability bias. Most easily missed is Old Ma's pitfall: the self giving a score, the self living, and the self remembering are not identical.

\section{Your Question: A Pill without Side Effects}
Return to the classroom.

Ms Gu's slips produced 6.8. Imagine visiting her ten years later and finding two new objects on her desk.

One is a pair of glasses. Wear them and any day acquires the texture of your happiest remembered day: warmer colours, kinder people, even school meals tasting of birthdays. Safe, no side effects, removable at any time.

The other is a pill. Take it and you will never again feel anxious or empty or lose sleep over anything. Life satisfaction stays at 9 until death. Again, safe and without side effects. Its only instruction says: you will achieve no more, love nobody more, and be loved by no additional people. You will simply feel wonderful, always.

Ms Gu asks: which would you take, or neither?

Answer with reasons. Before deciding, weigh again what this chapter has placed in your hands:

Bentham's ledger: feeling is the unavoidable judge, suffering the undeniable evil. For someone in pain, the pill may be mercy.

Yan Hui's unchanged joy: over two millennia ago, someone showed that joy could rest on deeper foundations than feeling.

Nozick's machine: most people hesitate before perfect feelings, revealing an attachment to reality and authorship.

Existentialism's reminder: meaning may grow precisely from anxiety, emptiness, and sleeplessness. Deleting them with one click may remove more than suffering.

Finally, remember Old Ma's question. The pill promises permanent satisfaction for ``you.'' Which you?

There is no standard answer. But from today, when someone asks whether you are happy, you can at least ask back: which kind do you mean?
